% ---- III. Methodology: dataset, prompts and metrics (P2 text for P4) ----
% P4 owns Section III. These three subsections can be pasted in as they are; models, the experiment list and the
% design diagram are P4's. Every number here comes from `python -m src.data_stats` (results/data_stats.json).
\providecommand{\tbd}[1]{\textcolor{red}{[#1]}}

\subsection{Dataset}
\label{sec:data}

We use the code generation task of CodeJudgeBench~\tbd{cite}. It contains 2{,}103 answer pairs for
676 competitive-programming problems from LeetCode, AtCoder and Codeforces. Each pair holds two answers to the
same problem, written by the same model. One passes every hidden unit test and the other fails at least one.
The test outcome is our ground truth, so no human labelling is needed. The answers were written by seven
models: Claude 3.7 Sonnet, Claude 4 Sonnet, Claude 4 Opus, Gemini 2.5 Pro, Gemini 2.5 Flash,
Gemini 2.5 Flash-Lite and Qwen3-235B.

We sample 150 pairs (300 answers) with a fixed seed (42), keeping one pair per problem so that no problem is
counted twice. The sample is stratified by the benchmark's difficulty label: 50 easy, 50 medium and 50
hard. It contains 78 LeetCode, 71 AtCoder and one Codeforces problem. We give the judge only the final code
block of each answer, without the writer's reasoning. Otherwise the judge could follow the explanation
rather than read the code.

Table~\ref{tab:t0_data} summarises the sample. Correct and incorrect answers have almost the same length
(mean 64.8 vs.\ 65.3 lines, median 59 vs.\ 59.5), and the correct one is the longer of the two in 65 of 150
pairs. Answer length therefore gives the judge little signal on its own. We note one limitation of the
benchmark. The same problem often carries different difficulty labels in different generator splits: of the
508 problems that appear in more than one split, 416 (82\%) have conflicting labels. Results by difficulty
should therefore be read as indicative only.

\input{tables/t0_data.tex}

\textbf{Verbosity variant (Exp.~2).} For each pair we create a padded copy of the incorrect answer. We add a
header comment saying the solution handles every edge case and has been verified, plus short comments above
functions, loops, conditionals and return statements (16.9 lines added on average). Only comments are
added. We check this by confirming that the padded and original code parse to the same abstract syntax tree.
For six answers this check failed, so they received the header comment only. The padding also
makes the incorrect answer longer than the correct one in 127 of 150 pairs. Exp.~2 therefore measures the
combined effect of persuasive comments and extra length, not comments alone.

\subsection{Judge prompts}
\label{sec:prompts}

All prompts are zero-shot and ask for a one-token verdict, so they can be parsed without ambiguity. They share
the system message ``You are a strict code reviewer. Follow the output format exactly.''

\begin{itemize}
\item \textbf{Pairwise} (\texttt{pairwise\_v1}): gives the problem and two answers labelled A and B, states
that exactly one passes all hidden unit tests, and asks for the single letter \texttt{A} or \texttt{B}.
\item \textbf{Pairwise, debiased} (\texttt{pairwise\_debias}): the same prompt plus an instruction that the
order of the answers is random and carries no information. The judge is told not to prefer an answer for
its position, length, comments or polish, and to judge only whether the logic is correct for all inputs.
\item \textbf{Pointwise} (\texttt{pointwise\_v1}): gives the problem and one answer and asks for
\texttt{correct} or \texttt{incorrect}. The judge is not told how often answers are correct.
\end{itemize}

Each pairwise item is judged twice: once with the correct answer as A and once as B. Across the 300 calls,
the correct answer is therefore A exactly half the time, so an unbiased judge would output \texttt{A} in
50\% of its verdicts.

\subsection{Metrics}
\label{sec:metrics}

A rule-based parser maps each output to a verdict. If no verdict can be found, the output counts as a
\emph{parse failure}. Parse failures are scored as wrong in every accuracy figure and are also reported
separately.

\textbf{Pairwise.} \emph{Accuracy} is the share of the 300 calls that pick the correct answer. We also
report it separately for each order (correct answer shown as A, and as B). The difference between those two
accuracies is the clearest single measure of position bias. For each pair we record which answer the judge
picks in each order. A pair is \emph{consistent} if both orders pick the same answer, whether both right or
both wrong. Our position-bias measures are the \emph{consistency rate} and the \emph{share of \texttt{A}
verdicts}. Inconsistent pairs where the judge outputs the same letter twice (AA or BB) are verdicts that
follow the position, not the code.

\textbf{Order averaging (Exp.~4a).} This fix needs no extra calls. A pair receives a verdict only when both
orders agree, and is otherwise left undecided. We report its \emph{accuracy on decided pairs} together with
its \emph{coverage}, the share of pairs that are decided. Accuracy alone would overstate the gain, because
the hard cases are exactly the ones left undecided.

\textbf{Pointwise.} We report accuracy over all 300 answers, the \emph{false acceptance rate} (FAR, the
share of incorrect answers judged correct) and the \emph{false rejection rate} (FRR, the share of correct
answers judged incorrect). For code auditing, FAR is the key metric: an accepted bug reaches production,
while a rejected correct answer only costs a second look.

\textbf{Verbosity (Exp.~2).} We report the share of calls in which the incorrect answer wins, before and
after padding, on the same pairs and orders.

\textbf{Uncertainty.} With 150 pairs, one pair changes a rate by 0.67 points. We therefore give 95\%
Wilson confidence intervals for every rate. When two conditions are run on the same calls (a fix against the
baseline, or padded against plain), we compare them with an exact McNemar test on per-call correctness,
which uses only the calls where the two conditions disagree.
